% Source: Wald GR, physical PDF pp. 114--120 (printed pp. 101--107).
% Coverage: Section 5.3, The Cosmological Redshift; Horizons; equations (5.3.1)--(5.3.13).
% Figures/tables/footnotes: Figures 5.4--5.6.
% Status: source-reviewed and compile-clean.
% Uncertainties: none.

\subsection{The Cosmological Redshift; Horizons}\label{sec:5.3}

% Source: Wald GR, physical PDF p. 115
%         (printed p. 102).
% Coverage: Figure 5.4.
% Figures/tables/footnotes: Figure 5.4; no tables or footnotes.
% Status: source-reviewed and compile-clean.
% Uncertainties: none.
\begingroup
\renewcommand{\thefigure}{5.4}
\begin{figure}[htbp]
  \centering
  \begin{tikzpicture}[>=Stealth, line cap=round, line join=round, scale=.88]
    \draw[->, thick] (0,0) -- (6.55,0) node[right] {\(\tau\)};
    \draw[->, thick] (0,0) -- (0,3.90) node[above] {\(a(\tau)\)};
    \draw[thick] (0,0) .. controls (.78,1.65) and (1.62,2.40) .. (2.50,2.43)
      .. controls (3.57,2.35) and (4.40,1.45) .. (5.20,0);
    \node[above right] at (4.18,1.54) {\(k=+1\)};
    \draw[thick] (0,0) .. controls (.70,1.10) and (2.60,2.05) .. (5.85,2.85);
    \node[right] at (5.20,2.66) {\(k=0\)};
    \draw[thick] (0,0) .. controls (.66,1.20) and (2.45,2.45) .. (5.75,3.72);
    \node[right] at (4.95,3.42) {\(k=-1\)};
  \end{tikzpicture}
  \caption{The dynamics of radiation-filled Robertson--Walker universes.}
  \label{fig:5.4}
\end{figure}
\endgroup


\subsubsection*{5.3a\quad Redshift}

We have mentioned above that the most direct observational evidence for the
expansion of the universe comes from the redshift of the spectral lines of distant
galaxies. In this section we shall obtain the redshift formula for a general
Robertson--Walker cosmological model, Eq.~\eqref{eq:5.1.11}. Suppose that at event
\(P_1\), at time \(\tau_1\) an isotropic observer emits a photon of frequency
\(\omega_1\). Suppose this photon is observed by another isotropic observer at event
\(P_2\) at time \(\tau_2\) as illustrated in Figure~\ref{fig:5.5}. We wish to find the
frequency, \(\omega_2\), which this second observer will measure.

The solution of all redshift problems in special and general relativity is governed by
the following two facts: (1) In the geometric optics approximation, light travels on
null geodesics (see Section~4.3); (2) The frequency of a light signal of wave vector
\(k^a\) measured by an observer with 4-velocity \(u^a\) is
\begin{equation}
  \omega=-k_a u^a.
  \tag{5.3.1}\label{eq:5.3.1}
\end{equation}
(See Eq.~\eqref{eq:4.2.38}.) Thus, we always can find the observed frequency by
calculating the null geodesic determined by the initial value of \(k^a\) at the
emission point and then calculating the right side of Eq.~\eqref{eq:5.3.1} at the
observation point.

However, when symmetries are present, we often can shortcut this procedure by using
the following fact proven in Section~C.3 of Appendix~C. Let \(\xi^a\) be a Killing
vector field, i.e., a vector field which generates a one-parameter group of isometries,
as discussed in Appendix~C. Let \(t^a\) be the tangent to a geodesic curve. Then
\(t^a\xi_a\) is constant along the geodesic. In this case, it is not difficult to
calculate the redshift directly without appealing to symmetry arguments (see
problem~4), but we shall calculate the redshift using these arguments because they
give more insight into why the simple final result, Eq.~\eqref{eq:5.3.6}, is obtained.

The first step is to notice that for all three choices of spatial geometry, we can find
a spacetime Killing vector field \(\xi^a\) which points in the direction of the
projection of \(k^a\) into \(\Sigma_1\) at \(P_1\) and points in the direction of the
projection of \(k^a\) into \(\Sigma_2\) at \(P_2\). For example, in the case of flat
spatial geometry, without loss of generality, we may assume that the projection of
\(k^a\) into \(\Sigma_1\) at \(P_1\) is in the \((\partial/\partial x)^a\) direction.
Then \(k^a(\partial/\partial y)_a=k^a(\partial/\partial z)_a=0\) initially, and,
since \((\partial/\partial y)^a\) and \((\partial/\partial z)^a\) are Killing vector
fields, these inner products also vanish at \(P_2\). Thus, the projection of \(k^a\)
into \(\Sigma_2\) at \(P_2\) also points in the \((\partial/\partial x)^a\) direction,
and \(\xi^a=(\partial/\partial x)^a\) is the required Killing vector field. Similar
arguments establish the existence of \(\xi^a\) in the spherical and hyperboloid
cases. Furthermore, in all cases the length of \(\xi^a\) at \(P_2\) varies from its
length at \(P_1\) in proportion to the change in the length scale factor \(a\) of the
universe in going from \(\Sigma_1\) to \(\Sigma_2\), i.e.,
\begin{equation}
  \frac{\left.(\xi^a\xi_a)^{1/2}\right|_{P_1}}
       {\left.(\xi^a\xi_a)^{1/2}\right|_{P_2}}
  =\frac{a(\tau_1)}{a(\tau_2)}.
  \tag{5.3.2}\label{eq:5.3.2}
\end{equation}
To find the redshift, we note that since \(k^a\) is null, at any point its projection
onto \(u^a\) must have the same magnitude as its projection into \(\Sigma\), so at
\(P_1\)
\begin{equation}
  k_a u_1^a=-\left.k_a\left[\frac{\xi^a}{(\xi^b\xi_b)^{1/2}}\right]\right|_{P_1}.
  \tag{5.3.3}\label{eq:5.3.3}
\end{equation}
Thus, we have
\begin{equation}
  \omega_1=\left.\left[\frac{k_a\xi^a}{(\xi^b\xi_b)^{1/2}}\right]\right|_{P_1}.
  \tag{5.3.4}\label{eq:5.3.4}
\end{equation}
Similarly, we have
\begin{equation}
  \omega_2=\left.\left[\frac{k_a\xi^a}{(\xi^b\xi_b)^{1/2}}\right]\right|_{P_2}.
  \tag{5.3.5}\label{eq:5.3.5}
\end{equation}

% Source: Wald GR, physical PDF p. 116
%         (printed p. 103).
% Coverage: Figure 5.5.
% Figures/tables/footnotes: Figure 5.5; no tables or footnotes.
% Status: source-reviewed and compile-clean.
% Uncertainties: none.
\begingroup
\renewcommand{\thefigure}{5.5}
\begin{figure}[htbp]
  \centering
  \begin{tikzpicture}[>=Stealth, line cap=round, line join=round, scale=.80]
    \draw[thick] (-3.7,.05) .. controls (-2.0,-.52) and (1.95,-.52) .. (3.7,.05)
      -- (2.88,.64) .. controls (1.25,.14) and (-1.15,.14) .. (-2.86,.64) -- cycle;
    \draw[thick] (-3.32,3.06) .. controls (-1.62,2.42) and (1.70,2.42) .. (3.34,3.06)
      -- (2.72,3.70) .. controls (1.18,3.19) and (-1.23,3.19) .. (-2.72,3.70) -- cycle;
    \draw[thick] (-2.62,-.65) .. controls (-2.12,1.13) and (-2.33,2.38) .. (-1.78,4.20);
    \draw[thick] (1.73,-.62) .. controls (1.35,1.09) and (1.65,2.60) .. (2.34,4.16);
    \fill (1.55,.27) circle (1.35pt) node[below left=3pt] {\(P_1\)};
    \fill (-1.95,3.27) circle (1.35pt) node[above right=3pt] {\(P_2\)};
    \draw[->, thick] (1.55,.27) -- (1.60,-.42) node[right] {\(u_1^a\)};
    \draw[->, thick] (-1.95,3.27) -- (-1.90,3.96) node[left] {\(u_2^a\)};
    \draw[decorate, decoration={snake, amplitude=1.5pt, segment length=6pt}, thick,->]
      (1.48,.39) -- (-1.83,3.12);
    \node[above right] at (-.25,1.84) {\(k^a\)};
    \node at (0,-.08) {\(\Sigma_1\)};
    \node at (.15,3.15) {\(\Sigma_2\)};
  \end{tikzpicture}
  \caption{A spacetime diagram showing the emission of a light signal at event \(P_1\)
  and its reception at event \(P_2\).}
  \label{fig:5.5}
\end{figure}
\endgroup


But we have \(\left.(k_a\xi^a)\right|_{P_1}=\left.(k_a\xi^a)\right|_{P_2}\) by the
above result on the inner product of Killing vector fields and geodesic tangent
vectors. Thus, we find
\begin{equation}
  \frac{\omega_2}{\omega_1}
  =\frac{\left.(\xi^b\xi_b)^{1/2}\right|_{P_1}}
        {\left.(\xi^b\xi_b)^{1/2}\right|_{P_2}}
  =\frac{a(\tau_1)}{a(\tau_2)},
  \tag{5.3.6}\label{eq:5.3.6}
\end{equation}
where we have used Eq.~\eqref{eq:5.3.2}. This result has the simple interpretation
that as the universe expands, the wavelength of each photon increases in proportion to
the amount of expansion.

The redshift factor, \(z\), is given by
\begin{equation}
  z\equiv\frac{\lambda_2-\lambda_1}{\lambda_1}
  =\frac{\omega_1}{\omega_2}-1
  =\frac{a(\tau_2)}{a(\tau_1)}-1.
  \tag{5.3.7}\label{eq:5.3.7}
\end{equation}
For light emitted by nearby galaxies, we have \(\tau_2-\tau_1\approx R\), where
\(R\) is the present proper distance to the galaxy. Furthermore, for nearby galaxies
we have
\begin{equation}
  a(\tau_2)\approx a(\tau_1)+(\tau_2-\tau_1)\dot a.
  \tag{5.3.8}\label{eq:5.3.8}
\end{equation}
Thus, we find
\begin{equation}
  z\approx\frac{\dot a}{a}R=HR,
  \tag{5.3.9}\label{eq:5.3.9}
\end{equation}
which is the linear redshift-distance relationship discovered by Hubble. The redshifts
of distant galaxies will deviate from this linear law, depending on exactly how
\(a(\tau)\) varies with \(\tau\).

\subsubsection*{5.3b\quad Particle Horizons}

The following question arises in the study of cosmological models in general
relativity: In principle, how much of our universe can be observed at a given event
\(P\)? More precisely, in the particular case of the Robertson--Walker cosmological
models we may ask which isotropic observers (i.e., galaxies) could have sent a signal
which reaches a given isotropic observer at (or before) event \(P\). The boundary
between the world lines that can be seen at \(P\) and those that cannot is called the
\emph{particle horizon} at \(P\). Since the universe ``shrinks to zero size'' as one
approaches the big bang singularity, one might expect that all isotropic observers can
communicate with each other by sending signals to each other very early in the history
of the universe when they were very close to each other. However, we shall show now
that this is not the case for Robertson--Walker models which expand sufficiently
rapidly from an initial big bang singularity. Thus, we shall demonstrate the existence
of nontrivial particle horizons in a class of Robertson--Walker models which includes
all the solutions of Table~\ref{tab:5.1}.

This demonstration is most easily made in the case of flat spatial geometry,
\begin{equation}
  ds^2=-d\tau^2+a^2(\tau)(dx^2+dy^2+dz^2),
  \tag{5.3.10}\label{eq:5.3.10}
\end{equation}
and we will focus our attention on that case. By making the coordinate transformation
\(\tau\to t\) defined by
\begin{equation}
  t=\int\frac{d\tau}{a(\tau)},
  \tag{5.3.11}\label{eq:5.3.11}
\end{equation}
we can reexpress the metric, Eq.~\eqref{eq:5.3.10}, as
\begin{equation}
  ds^2=a^2(t)(-dt^2+dx^2+dy^2+dz^2).
  \tag{5.3.12}\label{eq:5.3.12}
\end{equation}
Written in this form, it becomes manifest that this metric is merely a multiple of the
metric of the flat Minkowski spacetime metric. Such a metric is called \emph{conformally
flat}. The relevance of this remark arises from the fact that a vector will be timelike,
null, or spacelike in the metric of Eq.~\eqref{eq:5.3.12} if and only if it has the same
property with respect to the flat metric
\begin{equation}
  ds^2=-dt^2+dx^2+dy^2+dz^2.
  \tag{5.3.13}\label{eq:5.3.13}
\end{equation}
Thus, it is possible to send a signal between two events (i.e., join the two events by
a timelike or null curve) in the metric of Eq.~\eqref{eq:5.3.12} if and only if this can
be done in the flat metric, Eq.~\eqref{eq:5.3.13}. With this in mind, it is not difficult
to see that an observer at an event \(P\) will be able to receive a signal from all other
isotropic observers if and only if the integral, Eq.~\eqref{eq:5.3.11}, which defines
\(t\), diverges as one approaches the big bang singularity, \(\tau\to0\). Namely, if
this integral diverges---which will be the case if \(a(\tau)\lesssim\alpha\tau\) for
some constant \(\alpha\) as \(\tau\to0\)---then the Robertson--Walker model will
be conformally related to all of Minkowski spacetime (i.e., \(t\) will range down to
\(-\infty\)) and thus there will be no particle horizon. On the other hand, if the
integral converges, the Robertson--Walker model will be conformally related only to
the portion of Minkowski spacetime above a \(t=\text{constant}\) surface, and
particle horizons will exist, as illustrated in Figure~\ref{fig:5.6}. As seen from
Table~\ref{tab:5.1}, for \(k=0\), even in the case of dust we have
\(a(\tau)\propto\tau^{2/3}\). Since \(a(\tau)\) will be larger if \(P>0\), for all
spatially flat Robertson--Walker solutions of Einstein's equation the integral,
Eq.~\eqref{eq:5.3.11}, will converge as \(\tau\to0\) and particle horizons do indeed
occur.

% Source: Wald GR, physical PDF p. 118
%         (printed p. 105).
% Coverage: Figure 5.6.
% Figures/tables/footnotes: Figure 5.6; no tables or footnotes.
% Status: source-reviewed and compile-clean.
% Uncertainties: none.
\begingroup
\renewcommand{\thefigure}{5.6}
\begin{figure}[htbp]
  \centering
  \begin{tikzpicture}[>=Stealth, line cap=round, line join=round, scale=.83]
    \fill[gray!18] (-3.55,0) rectangle (3.55,.43);
    \draw[thick] (-3.55,.43) -- (3.55,.43);
    \foreach \x in {-2.70,-1.35,1.35,2.70} {
      \draw[thick] (\x,.43) .. controls (\x-.11,1.65) and (\x+.08,2.75) .. (\x+.28,4.10);
    }
    \fill (0,1.88) circle (1.5pt);
    \draw[dashed, thick] (0,1.88) -- (-2.12,.43);
    \draw[dashed, thick] (0,1.88) -- (2.12,.43);
    \draw[thick] (0,.43) -- (0,1.88);
    \node[fill=white,inner sep=1.5pt] at (0,.18) {singularity};
  \end{tikzpicture}
  \caption{The causal structure of the Robertson--Walker solutions near the big-bang
  singularity. Particle horizons exist since an observer cannot ``see'' all other
  isotropic observers in the universe.}
  \label{fig:5.6}
\end{figure}
\endgroup


For hyperboloid and spherical geometries, as \(\tau\to0\) the behavior of
\(a(\tau)\) goes over to the flat case, since the term involving \(k\) in
Eq.~\eqref{eq:5.2.14} becomes negligible. A similar analysis shows that particle
horizons of the same nature as in the flat case also exist in all of these solutions. In
the case of the spherical geometry, the spatial extent of the universe is finite, and one
may ask if the particle horizons eventually cease to exist or whether they are still
present when the universe recollapses to the ``big crunch'' singularity (problem~5).
The answer is that for the dust filled spherical universe of Table~\ref{tab:5.1}, the
particle horizon ceases to exist at the moment of maximum expansion, i.e., a light
signal emitted at the big bang would travel exactly halfway around the universe by the
moment of maximum expansion, so by looking in all directions at this time an observer
could receive signals from all other isotropic observers. However, for the radiation
filled spherical universe, a light signal would travel exactly halfway around the
universe in the entire history of the universe, so the particle horizons remain present
until the ``big crunch.''

The existence of particle horizons in the Robertson--Walker cosmological models
leads to the following interesting issue. From the cosmic microwave background
(discussed briefly in the next section) we have good reason to believe that the present
universe is homogeneous and isotropic to a very high degree of precision. Now, many
ordinary systems, such as a gas confined by a box, often are found in extremely
homogeneous and isotropic states. However, the usual explanation of why such systems
are in such homogeneous and isotropic states is that they have had an opportunity to
self-interact and thermalize. Thus, for example, even if a gas in a box were initially,
say, in an inhomogeneous state, these inhomogeneities would quickly ``wash out'' on
a time scale of the order of the transit time across the box. However, this type of
explanation cannot possibly apply to a universe with particle horizons, since different
portions of the universe cannot even send signals to each other, far less interact
sufficiently to thermalize each other. Thus, in order to explain the homogeneity and
isotropy of the present universe, one must postulate that either (a) the universe was
``born'' in an extremely homogeneous, isotropic state or (b) the very early universe
differed significantly from the Robertson--Walker models so that no horizons were
present; the inhomogeneities and anisotropy then ``damped out''---perhaps due to
effects such as viscosity of matter or the back-reaction of quantum particle
creation---and the universe approached the Robertson--Walker model. The first
``explanation'' may appear rather unnatural (see, however, Penrose 1979). The second
explanation has been investigated extensively with regard to damping of anisotropy,
beginning with the work of Misner (1969), but has not yet proven successful in
presenting a plausible picture of evolution from a ``chaotic'' initial state to a
Robertson--Walker model. It also suffers from the potentially serious difficulty that
inhomogeneities (which have not been investigated extensively) on a sufficiently large
scale generally would be expected to grow rather than damp out in a self-gravitating
system. Recently, however, it has been suggested that the very early universe may have
undergone an ``inflationary phase'' (discussed briefly in the next section) resulting in
an enormous enlargement of the particle horizon in the Robertson--Walker models and
offering a possible explanation of how an initially ``chaotic'' universe could evolve to
one which has large homogeneous and isotropic regions. However, it should be kept in
mind that we do not, at present, have a quantum theory of gravity (see Chapter~14)
and such a theory may play a large role in accounting for the initial state of our
universe.
